\chapter*{Prólogo}
\addcontentsline{toc}{chapter}{Prólogo}

La presente investigación surge de la observación directa de un problema
recurrente en el sector bancario peruano y latinoamericano: la toma de decisiones
tecnológicas sobre interfaces digitales se realiza con frecuencia de manera
informal, fragmentada y sin criterios que integren las particularidades del entorno
financiero.

Durante el ejercicio profesional en el ámbito del desarrollo de software para
entidades bancarias, fue posible constatar que la elección de tecnologías
de interfaz de usuario --decisión con impacto directo en la seguridad, el cumplimiento
normativo y la experiencia del cliente-- se apoyaba mayoritariamente en
preferencias individuales del equipo técnico o en tendencias del mercado, sin un
proceso estructurado que considerara el contexto institucional específico.

Esta realidad adquiere especial relevancia en un momento en que la
transformación digital ha dejado de ser una opción estratégica para convertirse en
una condición de supervivencia competitiva. Los frameworks Angular, React y
Vue.js concentran la gran mayoría de las implementaciones de canales digitales
bancarios a nivel global; sin embargo, la literatura especializada no ofrece una
metodología diseñada específicamente para evaluar y seleccionar entre ellos en
entornos financieros regulados. Los marcos de evaluación tecnológica existentes
son de propósito general y no contemplan dimensiones críticas para la banca,
como el cumplimiento de PCI-DSS, la integración con arquitecturas legacy de
larga data, o la resiliencia operacional exigida por regulaciones emergentes como
DORA.

El propósito de esta tesis es cubrir esa brecha metodológica, aportando a las
instituciones bancarias un instrumento sistemático, replicable y adaptable que les
permita tomar decisiones tecnológicas mejor fundamentadas. Se espera que los
resultados sean de utilidad no solo para la institución estudiada en el caso de
estudio, sino para cualquier entidad financiera que enfrente procesos similares de
modernización de sus canales digitales en el contexto latinoamericano.

Durante el desarrollo de esta investigación surgió un segundo factor que
obligó a revisar el planteamiento inicial: la rápida consolidación de
herramientas de generación de código asistida por Inteligencia Artificial.
Lo que en un primer momento se abordó como un criterio adicional entre
varios, terminó por convertirse en el eje que reorganiza el propio alcance
de la tesis. Si la IA generativa nivela la brecha de talento y curva de
aprendizaje entre frameworks --un factor tradicionalmente organizacional--
entonces lo que distingue a una tecnología de interfaz de usuario apta para
la banca deja de ser su ecosistema o su popularidad, y pasa a ser aquello
que ninguna IA puede sustituir: su solidez técnica y su compatibilidad
regulatoria. Ese giro, sugerido durante la revisión de esta tesis ante el
jurado, es el que da forma al título y al modelo final que aquí se
presenta.

\vspace{1cm}
\begin{flushright}
El autor
\end{flushright}
